\clearpage
\fancyhead[L]{\small GRAHAM'S PLANET.}
\SourceLeaf{482}
\begin{center}
\rule{0.55\linewidth}{0.4pt}\\[2pt]
Astronomische Nachrichten. Band XXVII. Nr. 639. Seite 235. 236. 1848 Juni 10.\\[2pt]
\rule{0.55\linewidth}{0.4pt}
\end{center}
\begin{flushright}Göttingen 1848 Mai\raisebox{.35ex}{,} 22.\end{flushright}

Den neuen Planeten habe ich gleich am Tage des Empfanges Ihrer Benachrichtigung am Meridiankreise beobachtet, und auch zwei Tage nachher noch einmal. Das erste mal war der Planet sehr lichtschwach, aber das zweite mal, wo die Luft mit Moorrauch erfüllt war, war diese Schwäche so gross, dass die Beobachtung nur wenig Vertrauen verdient. Indessen setze ich doch beide her.

\begin{center}
\textsc{Graham}'s Planet.\\[3pt]
\setlength{\tabcolsep}{4pt}
\begin{tabular}{l|r|r|r}
1848 & Göttinger m. Z. & AR. & Decl. \\\hline
Mai 9 & $11^{\rm u}\ 31^{\rm m}\ 3^{\rm s}95$ & $220^\circ\ 38'\ 6''6$ & $-11^\circ\ 57'\ 48''2$ \\
\hfill 11 & $11\ 21\ 14.7$ & $220\ 8\ 40.5$ & $-11\ 53\ 10.3$ \\
\end{tabular}
\end{center}

\clearpage
\hypertarget{native-482-june4}{}
\begin{center}
\rule{0.55\linewidth}{0.4pt}\\[2pt]
Astronomische Nachrichten. Band XXVII. Nr. 641. Seite 265. 266. 1848 Juni 29.\\[2pt]
\rule{0.55\linewidth}{0.4pt}
\end{center}
\begin{flushright}Göttingen 1848 Juni 4.\end{flushright}

Meinen beiden Ihnen am 22. Mai übersandten Beobachtungen des \textsc{Graham}schen Planeten kann ich nur noch eine, gleichfalls am Meridiankreise gemachte, beifügen:

\begin{center}
1848 Mai 24.\quad $10^{\rm u}\ 18^{\rm m}\ 39^{\rm s}0$ M. Z.\\
G. A. $217^\circ\ 15'\ 58''8$\qquad Abw. $11^\circ\ 29'\ 29''65$ Südl.
\end{center}

Der Planet war fast noch schwerer zu beobachten als am 11ten. Da jedoch nun der Zeitraum fast verdoppelt war, so habe ich den Herrn Professor \textsc{Goldschmidt} sogleich nachher zu einer neuen Berechnung der Bahn veranlasst, wovon folgende Elemente das Resultat sind:

\begin{center}
\textsc{Graham}'s Planet.\\[3pt]
\small
\setlength{\tabcolsep}{3pt}
\begin{tabular}{lr}
Epoche, Mittl. Länge 1848 Mai 1. $0^{\rm u}$ M. Z. Berl. \mbox{.\ .} & $215^\circ\ 23'\ 57''6$ \\
Perihel \mbox{.\ .\ .\ .\ .\ .\ .\ .\ .\ .\ .\ .\ .\ .\ .\ .\ .} & $71\ 41\ 23.9$ \\
Aufsteigender Knoten \mbox{.\ .\ .\ .\ .\ .\ .\ .\ .\ .\ .\ .} & $68\ 29\ 31.5$ \\
\multicolumn{2}{l}{\quad Alle drei vom mittl. Aequinoctium 1848 Mai 1 gezählt.} \\
Neigung der Bahn \mbox{.\ .\ .\ .\ .\ .\ .\ .\ .\ .\ .\ .\ .} & $5\ 35\ 23.2$ \\
Excentricitätswinkel \mbox{.\ .\ .\ .\ .\ .\ .\ .\ .\ .\ .\ .\ .} & $7\ 6\ 17.7$ \\
Logarithmus der halben grossen Axe \mbox{.\ .\ .\ .\ .\ .\ .} & $0.3776396$ \\
Mittlere tägliche siderische Bewegung \mbox{.\ .\ .\ .\ .\ .} & $962''8247$ \\
\end{tabular}
\end{center}

Gegründet waren diese Elemente auf die \textsc{Graham}'sche Beobachtung vom 26. April, die Berliner vom 11. Mai, und die meinige vom 24. Mai. Herr Prof. \textsc{Goldschmidt} hat diese Elemente mit allen mir bisher bekannt gewordenen Meridianbeobachtungen verglichen; ich lege diese Vergleichung auf einem besondern Blatt bei. In meinen Beobachtungen am Meridiankreise ist ein Stillstand eingetreten, da das Instrument einer daran vorzunehmenden Abänderung wegen jetzt abgenommen ist.
